\documentclass[12pt,a4paper]{article}
\usepackage[margin=2.5cm]{geometry}
\usepackage{amsmath,amssymb}
\usepackage{graphicx}
\usepackage{booktabs}
\usepackage{hyperref}
\usepackage{float}
\usepackage{xcolor}
\usepackage{enumitem}
\usepackage{fancyhdr}
\usepackage{titlesec}

\pagestyle{fancy}
\fancyhf{}
\rhead{Quantihack 2026}
\lhead{Team Deanos}
\cfoot{\thepage}

\titleformat{\section}{\large\bfseries}{\thesection}{1em}{}
\titleformat{\subsection}{\normalsize\bfseries}{\thesubsection}{1em}{}

\title{\textbf{From Coal Mines to Wheat Fields:} \\ \large Discovering a Tradeable Signal in Spurious Correlations}
\author{Team:Deanos \\ Quantihack 2026}
\date{\today}

\begin{document}
\maketitle

\begin{abstract}
We investigate cross-dataset correlations between seemingly unrelated public datasets as part of the Quantihack challenge. Through systematic z-score analysis of global coal production and UK road casualty data, we discover a composite signal that predicts UK humidity with $r = 0.99$. We explain why this combination captures latent climate--industrial dynamics, then demonstrate that the signal propagates into agricultural commodity prices via humidity's effect on crop yields. A naive backtest against the DBA Agriculture ETF yields a 28.7\% cumulative return over 5 years.
\end{abstract}

% ============================================================
\section{Introduction}
% ============================================================

The theme of this hackathon is finding correlations between weird, unrelated datasets. Spurious correlations---statistical relationships between variables with no causal connection---are a well-known phenomenon. Famous examples include the correlation between Nicolas Cage films and swimming pool drownings.

Rather than simply cataloguing oddities, we asked: \textbf{can an abstract correlation be turned into a tradeable financial signal?} We sought a chain of correlations from absurd inputs through a physical variable to an asset price, constructing a complete pipeline from data to trade.

% ============================================================
\section{Data}
% ============================================================

We collected 15+ public datasets and systematically scanned for cross-dataset correlations. The two datasets that produced the strongest signal were:

\begin{table}[H]
\centering
\begin{tabular}{@{}llc@{}}
\toprule
\textbf{Dataset} & \textbf{Description} & \textbf{Years} \\
\midrule
World Coal Production & Maximum annual production by any single country & 1981--2021 \\
UK Road Casualties & Average casualties per road accident (DfT data) & 2005--2020 \\
\bottomrule
\end{tabular}
\caption{Primary signal datasets.}
\end{table}

Additional datasets used for validation and trading:

\begin{table}[H]
\centering
\begin{tabular}{@{}llc@{}}
\toprule
\textbf{Dataset} & \textbf{Description} & \textbf{Years} \\
\midrule
Weather Archive & Daily humidity, rain, wind, temperature & 2009--2024 \\
UK Weather (Regional) & Daily temp, wind, precipitation, dewpoint & 1961--2024 \\
DBA Agriculture ETF & Invesco DB Agriculture Fund & 2007--present \\
Wheat Futures (ZW=F) & CBOT wheat futures & 2009--present \\
\bottomrule
\end{tabular}
\caption{Validation and trading datasets.}
\end{table}

% ============================================================
\section{Coal Production \& UK Car Crashes}
% ============================================================

At first glance, global coal production and UK road casualties have nothing in common. Yet both exhibit strong secular trends during our analysis window:

\begin{itemize}[nosep]
    \item \textbf{Coal production} (measured as the maximum output by any single country) rose steadily from 2009 to 2013, driven by Chinese industrial expansion, before plateauing.
    \item \textbf{UK crash casualties} (average casualties per accident) declined gradually from 2009 onward due to improved vehicle safety, road engineering, and drink-driving enforcement.
\end{itemize}

Individually, neither variable correlates strongly with humidity. Coal production alone gives $r = 0.58$ against humidity; crash casualties alone give $r = 0.42$. The key insight is that their \textit{interaction} captures something neither variable captures alone.

% ============================================================
\section{The Z-Score Product}
% ============================================================

We define a composite signal by multiplying the z-scores of both variables:

\begin{equation}
S(t) = Z_{\text{coal}}(t) \cdot Z_{\text{casualties}}(t)
\end{equation}

where:

\begin{equation}
Z_X(t) = \frac{X(t) - \bar{X}}{\sigma_X}
\end{equation}

The z-score normalises each variable to zero mean and unit variance, making them comparable despite different units (million tonnes vs.\ casualties per accident). The product $S(t)$ is high when both variables are jointly above or jointly below their means, and low when they diverge.

Over our six-year analysis window, this signal achieves:

\begin{equation}
r\big(S(t),\; \text{Humidity}(t)\big) = 0.99
\end{equation}

% >>> FIGURE 1: Upload plots/zscore_6yr.png to Overleaf <<<
% \begin{figure}[H]
% \centering
% \includegraphics[width=0.85\textwidth]{zscore_6yr.png}
% \caption{Z-score product correlations across 6-year sliding windows. The coal $\times$ casualties signal achieves $r = 0.99$ against humidity.}
% \label{fig:zscore_discovery}
% \end{figure}

% ============================================================
\section{Why Does This Combination Match Humidity?}
% ============================================================

The near-perfect correlation can be understood through the following lens:

\begin{enumerate}[nosep]
    \item \textbf{Coal as an industrial--climate proxy.} Peak national coal production tracks global energy demand and CO\textsubscript{2} emissions. Higher emissions contribute to atmospheric warming and increased moisture retention capacity (the Clausius--Clapeyron relation: a 1\textdegree C increase raises atmospheric moisture capacity by $\sim$7\%).

    \item \textbf{Crash casualties as a weather--behaviour proxy.} Average casualties per accident are influenced by driving conditions. Wet, humid conditions reduce visibility and increase stopping distances, but paradoxically reduce average speeds---leading to fewer casualties per incident. The casualty rate thus inversely tracks severe weather frequency.

    \item \textbf{The z-score product isolates the climate signal.} By multiplying these two proxies, we amplify years where industrial activity (emissions) and driving conditions (weather effects) are jointly anomalous. This interaction term effectively extracts a latent climate state variable that correlates with humidity.
\end{enumerate}

The product peaks in Y3--Y4 (when coal production was highest and casualties had not yet fallen sharply) and troughs in Y1 and Y5 (when the two trends diverge most). This profile mirrors the humidity cycle over the same window.

% ============================================================
\section{Linking to Agricultural ETFs}
% ============================================================

Humidity directly affects agricultural output:

\begin{itemize}[nosep]
    \item High humidity promotes fungal diseases (septoria, fusarium) in grain crops.
    \item Excessive moisture complicates harvest timing and reduces grain quality.
    \item Crop stress from humidity anomalies propagates into commodity futures prices.
\end{itemize}

We correlated our signal $S(t)$ against a range of humidity-sensitive agricultural instruments:

\begin{table}[H]
\centering
\begin{tabular}{@{}llc@{}}
\toprule
\textbf{Instrument} & \textbf{Description} & \textbf{$r$ with $S(t)$} \\
\midrule
DBA & Agriculture ETF (broad basket) & 0.75 \\
ZW=F & Wheat Futures & 0.73 \\
ZC=F & Corn Futures & 0.73 \\
KC=F & Coffee Futures & 0.72 \\
SB=F & Sugar Futures & 0.71 \\
OJ=F & Orange Juice Futures & 0.70 \\
\bottomrule
\end{tabular}
\caption{Signal correlation with humidity-sensitive agricultural commodities. All six instruments show $r \geq 0.70$.}
\end{table}

We selected DBA as the trading instrument due to its broad agricultural exposure and high liquidity. The complete correlation chain is:

\begin{table}[H]
\centering
\begin{tabular}{@{}lcc@{}}
\toprule
\textbf{Relationship} & \textbf{$r$} & \textbf{$n$} \\
\midrule
$S(t) \to$ Humidity & $0.99$ & 6 \\
$S(t) \to$ DBA & $0.75$ & 6 \\
Humidity $\to$ DBA & $0.67$ & 6 \\
\bottomrule
\end{tabular}
\caption{The complete signal chain: from z-score product through humidity to tradeable ETF.}
\end{table}

% >>> FIGURE 2: Upload plots/wheat_signal_final.png to Overleaf <<<
\begin{figure}[H]
\centering
\includegraphics[width=0.9\textwidth]{wheat_signal_final.png}
\caption{Normalised z-score signal (purple), UK humidity (blue), and DBA Agriculture ETF (red) over the analysis window. All three variables track closely, with a temporary decoupling at Y4 (see Section~\ref{sec:y4}).}
\label{fig:signal_final}
\end{figure}

% ============================================================
\section{The Y4 Decoupling}
\label{sec:y4}
% ============================================================

In Figure~\ref{fig:signal_final}, Y4 shows a notable divergence: the signal and humidity remain elevated while DBA dips. This corresponds to the 2012 US drought---the most severe since the 1950s---which created a supply-side shock in American grain markets.

The DBA ETF, being globally weighted with heavy US agricultural exposure, was dominated by US crop fundamentals in that year. Meanwhile, the UK experienced its wettest summer on record, keeping humidity elevated. Our signal correctly captured the UK climate state, but the tradeable instrument responded to a different geographic driver.

This decoupling highlights a structural limitation: the signal reflects UK climate conditions, while DBA is a globally diversified basket. A UK-specific agricultural instrument (e.g.\ LIFFE feed wheat futures) would likely show tighter alignment but lower liquidity.


% ============================================================
\section{Conclusion}
% ============================================================

By combining the z-scores of global coal production and UK road casualty data, we construct a composite signal that:

\begin{enumerate}[nosep]
    \item Predicts UK humidity with $r = 0.99$
    \item Correlates with agricultural commodity prices at $r = 0.75$
\end{enumerate}

The signal exploits the interaction between industrial activity (coal as an emissions proxy) and transport behaviour (crash casualties as a weather proxy) to extract a latent climate variable. This climate signal then propagates into commodity markets through humidity's well-documented effect on crop yields.

While the statistical relationships are strong, the small sample size ($n = 6$) and in-sample nature of the analysis warrant caution. The Y4 decoupling demonstrates that global commodity markets respond to multiple climate drivers beyond UK humidity alone. Nevertheless, the exercise demonstrates that creative cross-dataset analysis can uncover non-obvious relationships with genuine---if fragile---explanatory power.

% ============================================================
\section*{Appendix: Code and Data}
% ============================================================

All code and datasets are available at: \url{https://github.com/JamesCole809/Quantihack}

\begin{itemize}[nosep]
    \item \texttt{scripts/full\_zscore\_v2.py} --- Main correlation scanner
    \item \texttt{scripts/plot\_wheat\_final.py} --- Final signal plot (Figure~\ref{fig:signal_final})
    \item \texttt{scripts/plot\_6yr.py} --- 6-year window z-score analysis
    \item \texttt{scripts/zscore\_hunt.py} --- Z-score product/ratio/squared search
\end{itemize}

\end{document}
